\documentclass[11pt,a4paper]{article}
\usepackage[margin=1.8cm]{geometry}
\usepackage{amsmath,amssymb,array,tabularx,fontspec,cancel,tikz,graphicx}
\setmainfont{FreeSerif}
\setlength{\parindent}{0pt}
\setlength{\parskip}{0pt}
% \lqn{number}{text}: long question with its printed number
\newcommand{\lqn}[2]{\par\medskip\noindent\hangindent=3em\hangafter=1
  \makebox[3em][l]{\textbf{#1.}}#2\par}
\newcommand{\tfig}[1]{\par\smallskip\noindent\hspace*{3em}\begin{tikzpicture}[scale=0.8,font=\small]#1\end{tikzpicture}\par}
\newcommand{\pfig}[2]{\par\smallskip\noindent\hspace*{3em}\includegraphics[width=#2]{#1}\par}
% data table helper: \dtab{cols}{row1}{row2...} -> centered-ish bordered tabular indented
\newcommand{\dtab}[2]{\par\smallskip\noindent\hspace*{3em}\begin{tabular}{#1}\hline #2\\ \hline\end{tabular}\par\smallskip}
\newcommand{\nums}[1]{\par\smallskip\noindent\hspace*{3em}\parbox{0.9\linewidth}{#1}\par\smallskip}

\begin{document}
\begin{center}
  {\Huge\bfseries The Hope Science Academy}\\[3pt]
  {\small Excellence is our standard not goal}\\[2pt]
  {\small 247-E-1, (Opposite LDA School) Johar Town, Lahore. PH: 0300-8194789, 0322-4157001}
\end{center}
\medskip\noindent
\begin{tabularx}{\linewidth}{|X|X|X|X|}
\hline
\textbf{Student Name:} & \textbf{Roll Num:} & \textbf{Class:} 9th & \textbf{Paper Code:} \\ \hline
\textbf{Subject Name:} Mathematics. & \textbf{Time Allowed:} & \textbf{Total Marks:} 2653 & \textbf{Exam Date:} 09-Oct-2026 \\ \hline
\multicolumn{4}{|l|}{\textbf{Exam Syllabus:} FULL BOOK} \\ \hline
\end{tabularx}
\bigskip
\noindent\textbf{Q4. Write detailed answers of the following questions.}\hfill\textbf{1$\times$448=448}

\lqn{1}{If $a=\dfrac{2}{3},\ b=\dfrac{3}{2},\ c=\dfrac{5}{3}$ then verify the distributive properties over addition.}
\lqn{2}{Simplify the following: $\dfrac{5.(25)^{n+1}-25.(5)^{2n}}{5.(5)^{2n+3}-(25)^{n+1}}$}
\lqn{3}{Simplify the following: $\dfrac{3^{n}\times9^{n+1}}{3^{n-1}\times9^{n-1}}$}
\lqn{4}{Simplify the following: $\dfrac{(16)^{x+1}+20(4^{2x})}{2^{x-3}\times8^{x+2}}$}
\lqn{5}{Simplify the following: $\dfrac{5^{n+3}-6.5^{n+1}}{9\times5^{n}-2^{n}\times5^{n}}$}
\lqn{6}{If $x=3+\sqrt{8}$ then find the value of: $x^{2}-\dfrac{1}{x^{2}}$}
\lqn{7}{Find the rational numbers p and q such that $\dfrac{8-3\sqrt{2}}{4+3\sqrt{2}}=p+q\sqrt{2}$.}
\lqn{8}{If $x=3+\sqrt{8}$ then find the value of: $x^{4}+\dfrac{1}{x^{4}}$}
\lqn{9}{Simplify the following: $\dfrac{(25)^{\frac{3}{2}}\times(243)^{\frac{3}{5}}}{(16)^{\frac{5}{4}}\times(8)^{\frac{4}{3}}}$}
\lqn{10}{Simplify the following: $\dfrac{54\times\sqrt[3]{(27)^{2x}}}{9^{x+1}+216\,(3^{2x-1})}$}
\lqn{11}{Simplify the following: $\left(a^{\frac{1}{3}}+b^{\frac{2}{3}}\right)\times\left(a^{\frac{2}{3}}-a^{\frac{1}{3}}b^{\frac{2}{3}}+b^{\frac{4}{3}}\right)$}
\lqn{12}{Simplify the following: $\sqrt{\dfrac{(216)^{\frac{2}{3}}\times(25)^{\frac{1}{2}}}{(0.04)^{\frac{-3}{2}}}}$}
\lqn{13}{The diagram shows right angled $\triangle ABC$ in which the length of $\overline{AC}$ is $(\sqrt{3}+\sqrt{5})$cm. The area of $\triangle ABC$ is $(1+\sqrt{15})\,cm^{2}$. Find the length $\overline{AB}$ in the form $(a\sqrt{3}+b\sqrt{5})\,cm$, where a and b are integers.}
\tfig{\coordinate(A)at(0,0);\coordinate(B)at(0,2);\coordinate(C)at(4,0);
\draw(A)--(B)--(C)--cycle;\draw(0,.25)--(.25,.25)--(.25,0);
\node[left]at(A){$A$};\node[above]at(B){$B$};\node[right]at(C){$C$};
\node[below]at(2,0){$(\sqrt{3}+\sqrt{5})\,cm$};}
\lqn{14}{A rectangle has sides of length $2+\sqrt{18}\,m$ and $\left(5-\dfrac{4}{\sqrt{2}}\right)$m. Express the area of the rectangle in the form $a+b\sqrt{2}$, where a and b are integers.}
\lqn{15}{Find two numbers whose sum is 68 and difference is 22.}
\lqn{16}{The weather in Lahore was unusually warm during the summer of 2024. The TV news reported temperature as high as 48$^{\circ}$C. By using the formula, $\left({}^{o}F=\dfrac{9}{5}{}^{o}C+32\right)$ find the temperature as Fahrenheit scale.}
\lqn{17}{The sum of the ages of the father and son is 72 years. Six years ago, the father's age was 2 times the age of the son. What was son's age six years ago?}
\lqn{18}{Mirha sells a toy for Rs. 1520. What will the selling price be to get a 15\% profit?}
\lqn{19}{The annual income of Tayyab is Rs. 9,60,000, while the exempted amount is Rs. 1,30,000. How much tax would he have to pay at the rate of 0.75\%?}
\lqn{20}{Find the compound markup on Rs. 3,75,000 for one year at the rate of 14\% compounded annually.}
\lqn{21}{The sum of three consecutive odd integers is 51. Find the three integers.}
\lqn{22}{Abdullah picked up 96 balls and placed them into two buckets. One bucket has twenty-eight more balls than the other bucket. How many balls were in each bucket?}
\lqn{23}{Salma invested Rs. 3,50,000 in a bank, which paid simple profit at the rate of $7\dfrac{1}{4}\%$ per annum. After 2 years, the rate was increased to 8\% per annum. Find the amount she had at the end of 7 years.}
\lqn{24}{Hamail purchased a bicycle for Rs. 6590 and sold it for Rs.6850. Find the profit percentage.}
\lqn{25}{Umair bought a book for Rs. 850 and sold it for Rs. 720. What was his loss percentage?}
\lqn{26}{Saleem, Nadeem, and Tanveer earned a profit of Rs. 4,50,000 from a business. If their investments in the business are in the ratio 4: 7: 14, find each person's profit.}
\lqn{27}{If the simple profit on Rs. 6400 for 12 years is Rs. 3840. Find the rate of profit.}
\lqn{28}{Find the value of x in the following: $\log_{4}x=\dfrac{3}{2}$}
\lqn{29}{Find the value of x in the following: $\log_{2}1024=x$}
\lqn{30}{Find the value of x: log x = -3.434}
\lqn{31}{Find the value of x: log x = -1.5726}
\lqn{32}{Find the value of x: log x = 4.3561}
\lqn{33}{Find the value of x: log x = 2.0184}
\lqn{34}{Find the value of x in the following: log x = -2.1234}
\lqn{35}{Without using calculator, evaluate the following: $\dfrac{1}{3}\log_{3}8-\log_{3}18$}
\lqn{36}{Without using calculator, evaluate the following: $\dfrac{1}{3}\log_{4}64+2\log_{5}25$}
\lqn{37}{Without using calculator, evaluate the following: $\log_{3}12+\log_{3}0.25$}
\lqn{38}{Find the value of x in the following equation: $(81)^{x}=(243)^{x+2}$}
\lqn{39}{Find the value of x in the following equation: $\left(\dfrac{1}{27}\right)^{x-6}=27$}
\lqn{40}{Find the value of x in the following equation: $\log_{2}(x+1)-\log_{2}(x-4)=2$}
\lqn{41}{Find the values of the following with the help of logarithm table: $\dfrac{3.68\times4.21}{5.234}$}
\lqn{42}{Find the values of the following with the help of logarithm table: $4.67\times2.11\times2.397$}
\lqn{43}{Find the values of the following with the help of logarithm table: $\dfrac{(20.46)^{2}\times(2.4122)}{754.3}$}
\lqn{44}{Find the values of the following with the help of logarithm table: $\dfrac{\sqrt[3]{9.364}\times21.64}{3.21}$}
\lqn{45}{The formula to measure the magnitude of earthquakes is given by $M=\log_{10}\left(\dfrac{A}{A_{o}}\right)$. If amplitude (A) is 10,000 and reference amplitude (A$_{o}$) is 10. What is the magnitude of the earthquake?}
\lqn{46}{Abdullah invested Rs. 100,000 in a saving scheme and gains interest at the rate of 5\% per annum so that the total value of this investment after t years us Rs y. This is modelled by an equation y= 100,000 (1.05)$^{t}$, $t\geq0$. Find after how many years the investment will be double.}
\lqn{47}{Huria is hiking up a mountain where the temperature (T) decreases by 3\% (or a factor of 0.97) for every 100 metres gained in altitude. The initial temperature (T$_{i}$) at sea level is 20$^{o}$C. Using the formula $T=T_{i}\times0.97^{\frac{h}{100}}$, calculate the temperature at an altitude (h) of 500 metres.}
\lqn{48}{The decibel scale measures sound intensity using the formula $L=40\log_{10}\left(\dfrac{I}{I_{o}}\right)$. If a sound has an intensity (I) of $10^{6}$ times the reference intensity (I$_{o}$). What is the sound level in decibels?}
\lqn{49}{Find value of x in the following: $\left(\dfrac{1}{32}\right)^{2x}=64$}
\lqn{50}{Find the values of the following with the help of logarithm table: 319.8 x 3.543}
\lqn{51}{Find the values of the following with the help of logarithm table: $\dfrac{36.12\times750.9}{113.2\times9.98}$}
\lqn{52}{In the year 2016, the population of the city was 22 millions and was growing at a rate of 2.5 \% per year. The function p(t) = 22(1.025)$^{t}$ gives the population in millions, t years after 2016. Use the model to determine in which year the population will reach 35 millions. Round the answer to the nearest year.}
\lqn{53}{Consider the universal set $U=\{x: x \text{ is multiple of } 2 \text{ and } 0<x\leq30\}$, A=\{x;x is a multiple of 6\} and B=\{x:x is a multiple of 8\}: Draw a Venn diagram}
\lqn{54}{Let, $U=\{x: x \text{ is an integer and } 0<x\leq150\}$, $G=\{x: x=2^{m} \text{ for integer m and } 0\leq m\leq12\}$ and $H=\{x: x \text{ is a square}\}$: Find $G\cup H$}
\lqn{55}{Let, $U=\{x: x \text{ is an integer and } 0<x\leq150\}$, $G=\{x: x=2^{m} \text{ for integer m and } 0\leq m\leq12\}$ and $H=\{x: x \text{ is a square}\}$: Find $G\cup H$}
\lqn{56}{Let, $U=\{x: x \text{ is an integer and } 0<x\leq150\}$, $G=\{x: x=2^{m} \text{ for integer m and } 0\leq m\leq12\}$ and $H=\{x: x \text{ is a square}\}$: Find $G\cap H$}
\lqn{57}{Consider the sets $P=\{x; x \text{ is a prime number and } 0<x\leq20\}$ and $Q=\{x: x \text{ is advisor of } 210 \text{ and } 0<x\leq20\}$: Find $P\cap Q$}
\lqn{58}{Consider the sets $P=\{x; x \text{ is a prime number and } 0<x\leq20\}$ and $Q=\{x: x \text{ is advisor of } 210 \text{ and } 0<x\leq20\}$: Find $P\cup Q$}
\lqn{59}{If U=\{1,2,3,...,20\} and A=\{1,3,5,...,19\}, verify the following: $A\cup A'=U$}
\lqn{60}{If U=\{1,2,3,...,20\} and A=\{1,3,5,...,19\}, verify the following: $A\cap A'=\phi$}
\lqn{61}{In a class of 55 students, 34 like to play cricket and 30 like to play hockey. Also each student likes to play at least one of the two games. How many students like to play both games?}
\lqn{62}{In a group of 500 employees, 250 can speak Urdu, 150 can speak English, 50 can speak Punjabi, 40 can speak Urdu and English, 30 can speak both English and Punjabi, and 10 can speak Urdu and Punjabi. How many can speak all three languages?}
\lqn{63}{In sports events, 19 people wear blue shirts, 15 wear green shirts, 3 wear blue and green shirts, 4 wear a cap and green shirts. The total number of people with either a blue or green shirt or cap is 25. How many people are wearing caps?}
\lqn{64}{In a training session, 17 participants have laptops, 11 have tablets, 9 have laptops and tablets, 6 have laptops and books, and 4 have both tablets and books. Eight participants have all three items. The total number of participants with laptops, tablets, or books is 35. How many participants have books?}
\lqn{65}{A shopping mall has 150 employees labelled 1 to 150, representing the Universal set U. The employees fall into the following categories: Set A: 40 employees with a salary range of 30k-45k, labelled from 50 to 89. Set B: 50 employees with a salary range of 50k-80k, labelled from 101 to 150. Set C: 60 employees with a salary range of 100k-150k, labelled from 1 to 49 and 90 to 100. Find $(A'\cup B')\cap C$}
\lqn{66}{A shopping mall has 150 employees labelled 1 to 150, representing the Universal set U. The employees fall into the following categories: Set A: 40 employees with a salary range of 30k-45k, labelled from 50 to 89. Set B: 50 employees with a salary range of 50k-80k, labelled from 101 to 150. Set C: 60 employees with a salary range of 100k-150k, labelled from 1 to 49 and 90 to 100. Find $n\{A\cap(B^{c}\cap C^{c})\}$}
\lqn{67}{In a secondary school with 125 students participants in at least one of the following sports: cricket, football, or hockey: 60 students play cricket. 70 students play football, 40 students play hockey. 25 student play both cricket and football. 15 students play both football and hockey. 10 students play both cricket and hokey. How many students play all three sports?}
\lqn{68}{In a secondary school with 125 students participants in at least one of the following sports: cricket, football, or hockey: 60 students play cricket. 70 students play football, 40 students play hockey. 25 student play both cricket and football. 15 students play both football and hockey. 10 students play both cricket and hokey. Draw a Venn diagram showing the distribution of sports participants in all the games.}
\lqn{69}{A survey was conducted in which 130 people were asked about their favourite foods. The survey results showed the following information. 40 people said they liked nihari. 65 people said they liked biryani. 50 people said they liked korma. 20 people said they liked nihari and biryani. 35 people said they liked biryani and korma. 27 people said they liked nihari and korma. 12 people said they liked all three foods nihari, biryani, and korma. At least how many people like nihari, biryani or korma?}
\lqn{70}{A survey was conducted in which 130 people were asked about their favourite foods. The survey results showed the following information. 40 people said they liked nihari. 65 people said they liked biryani. 50 people said they liked korma. 20 people said they liked nihari and biryani. 35 people said they liked biryani and korma. 27 people said they liked nihari and korma. 12 people said they liked all three foods nihari, biryani, and korma. How many people did not like nihari, biryani or korma?}
\lqn{71}{A survey was conducted in which 130 people were asked about their favourite foods. The survey results showed the following information. 40 people said they liked nihari. 65 people said they liked biryani. 50 people said they liked korma. 20 people said they liked nihari and biryani. 35 people said they liked biryani and korma. 27 people said they liked nihari and korma. 12 people said they liked all three foods nihari, biryani, and korma. How many people like only one of the following foods, nihari, biryani, or korma?}
\lqn{72}{A survey was conducted in which 130 people were asked about their favourite foods. The survey results showed the following information. 40 people said they liked nihari. 65 people said they liked biryani. 50 people said they liked korma. 20 people said they liked nihari and biryani. 35 people said they liked biryani and korma. 27 people said they liked nihari and korma. 12 people said they liked all three foods nihari, biryani, and korma. Draw a Venn diagram.}
\lqn{73}{Consider the adjacent Venn diagram illustrating two non-empty sets, A and B. Determine the number of elements common to sets A and B .}
\pfig{fig73}{5.5cm}
\lqn{74}{Consider the adjacent Venn diagram illustrating two non-empty sets, A and B. Identify all the elements exclusively in set B and not in set A}
\pfig{fig74}{5.5cm}
\lqn{75}{Consider the adjacent Venn diagram illustrating two non-empty sets, A and B. Calculate union sets of A and B}
\pfig{fig75}{5.5cm}
\lqn{76}{Consider the adjacent Venn diagram representing the students enrolled in different courses in an IT institution: U = \{Students enrolled in IT institutions\}, A=\{Student enrolled in an Applied Robotics\}, G = \{Students enrolled in a Game Development\}, W = \{Students enrolled in a Web Designing\}. How many student enrolled in the applied Robotics course?}
\pfig{fig76}{5cm}
\lqn{77}{Consider the adjacent Venn diagram representing the students enrolled in different courses in an IT institution: U = \{Students enrolled in IT institutions\}, A=\{Student enrolled in an Applied Robotics\}, G = \{Students enrolled in a Game Development\}, W = \{Students enrolled in a Web Designing\}. Determine the total number of students enrolled in a Game Development.}
\pfig{fig76}{5cm}
\lqn{78}{Consider the adjacent Venn diagram representing the students enrolled in different courses in an IT institution: U = \{Students enrolled in IT institutions\}, A=\{Student enrolled in an Applied Robotics\}, G = \{Students enrolled in a Game Development\}, W = \{Students enrolled in a Web Designing\}. How many students are enrolled in the Game development and Web designing course?}
\pfig{fig76}{5cm}
\lqn{79}{Consider the adjacent Venn diagram representing the students enrolled in different courses in an IT institution: U = \{Students enrolled in IT institutions\}, A=\{Student enrolled in an Applied Robotics\}, G = \{Students enrolled in a Game Development\}, W = \{Students enrolled in a Web Designing\}. Identify the students enrolled in Web development but not Applied Robotics?}
\pfig{fig76}{5cm}
\lqn{80}{Consider the adjacent Venn diagram representing the students enrolled in different courses in an IT institution: U = \{Students enrolled in IT institutions\}, A=\{Student enrolled in an Applied Robotics\}, G = \{Students enrolled in a Game Development\}, W = \{Students enrolled in a Web Designing\}. How many student are enrolled in IT institutions?}
\pfig{fig76}{5cm}
\lqn{81}{Consider the adjacent Venn diagram representing the students enrolled in different courses in an IT institution: U = \{Students enrolled in IT institutions\}, A=\{Student enrolled in an Applied Robotics\}, G = \{Students enrolled in a Game Development\}, W = \{Students enrolled in a Web Designing\}. How many student enrolled in all three courses?}
\pfig{fig76}{5cm}
\lqn{82}{There are 95 secondary school students in a sports club. 58 students join the swimming club, and 50 join the tug-of-war club. How many students participated in both games?}
\lqn{83}{Mr. Saleem, a school teacher, has a small library in his house containing 150 books. He has two main categories for these books: Islamic and science. He categorized 70 books as Islamic books and 90 books as science books. There are 15 books that neither belong to the Islamic nor science books category. How many books are classified under both the Islamic and science categories?}
\lqn{84}{In a college, 45 teacher teach mathematics or physics or chemistry. Here is information about teachers who teach different subjects: 18 teach mathematics, 12 teach physics, 8 teach chemistry, 6 teach both mathematics and physics, 4 teach both physics and chemistry, 2 teach both mathematics and chemistry, How many teachers teach all three subjects?}
\lqn{85}{A survey of 130 customers in a shopping mall was conducted in which they were asked about buying preferences. The survey result showed the following statistics. 57 customers bought garments, 50 customers bought cosmetics, 46 customers bought electronics, 31 customers purchased both garments and cosmetics. 25 customers purchased both garments and electronics, 21 customers purchased both cosmetics and electronics, 12 customers purchased all three products i.e garments, cosmetics and electronics. How many of the customers bought at least one of the products, and electronics.}
\lqn{86}{A survey of 130 customers in a shopping mall was conducted in which they were asked about buying preferences. The survey result showed the following statistics. 57 customers bought garments, 50 customers bought cosmetics, 46 customers bought electronics, 31 customers purchased both garments and cosmetics. 25 customers purchased both garments and electronics, 21 customers purchased both cosmetics and electronics, 12 customers purchased all three products i.e garments, cosmetics and electronics. How many of the customers bought only one of the products: garments, cosmetics and electronics?}
\lqn{87}{A survey of 130 customers in a shopping mall was conducted in which they were asked about buying preferences. The survey result showed the following statistics. 57 customers bought garments, 50 customers bought cosmetics, 46 customers bought electronics, 31 customers purchased both garments and cosmetics. 25 customers purchased both garments and electronics, 21 customers purchased both cosmetics and electronics, 12 customers purchased all three products i.e garments, cosmetics and electronics. How many customers did not buy any of the three products?}
\lqn{88}{Given that $f(x)=ax+b+1$, where a and b are constant numbers. If $f(3)=8$ and $f(6)=14$, then find the value of a and b.}
\lqn{89}{Given that $g(x)=ax+b+5$, where a and b are constant numbers. If $g(-1)=0$ and $g(2)=10$, find the values of a and b.}
\lqn{90}{Consider the function $f(x)=cx^{2}+d$, where c and d are constant numbers. If $f(1)=6$ and $f(-2)=10$, then find the values of c and d.}
\lqn{91}{Let $c_{1},c_{2},c_{3}$ be three children and $m_{1},m_{2}$ be two men such that the father of both $c_{1},c_{2}$ is m$_{1}$ and father of $c_{3}$ is m$_{2}$. Find the relation \{(child,father)\}}
\lqn{92}{Consider $f(x)=ax+b+3$, where a and b are constant numbers. If $f(1)=4$ and $f(5)=9$, then find the value of a and b.}
\lqn{93}{Use Venn diagram to verify the following: $A-B=A\cup B'$}
\lqn{94}{Use Venn diagram to verify the following: $(A-B)'\cap B=B$}
\lqn{95}{Verify the properties for the sets, A,B and C given below: (i) Associativity of Union (ii) Associativity of intersection (iii) Distributivity of Union over intersection. (iv) Distributivity of intersection over union. $A=\{1,2,3,4\},\ B=\{3,4,5,6,7,8\},\ C=\{5,6,7,9,10\}$}
\lqn{96}{Verify the properties for the sets, A,B and C given below: (i) Associativity of Union (ii) Associativity of intersection (iii) Distributivity of Union over intersection. (iv) Distributivity of intersection over union. $A=\phi,\ B=\{0\},\ C=\{0,1,2\}$}
\lqn{97}{Verify the properties for the sets, A,B and C given below: (i) Associativity of Union (ii) Associativity of intersection (iii) Distributivity of Union over intersection. (iv) Distributivity of intersection over union. $A=N,\ B=Z,\ C=Q$}
\lqn{98}{Verify the Morgan's Laws for the following set: $U=\{1,2,3,\ldots,20\},\ A=\{2,4,6,\ldots,20\}$ and $B=\{1,3,5,\ldots,19\}$}
\lqn{99}{From suitable properties of union and intersection, deduce the following result: $A\cap(A\cup B)=A\cup(A\cap B)$}
\lqn{100}{From suitable properties of union and intersection, deduce the following result: $A\cup(A\cap B)=A\cap(A\cup B)$}
\lqn{101}{Given that $f(x)=ax+b$, where a and b are constant numbers, if $f(-2)=3$ and $f(4)=10$, then find the values of a and b.}
\lqn{102}{Consider the function $g(x)=mx^{2}+n$, where m and n are constant numbers, If g(4)=20 and g(0)=5, find the values of m and n.}
\lqn{103}{Out of the 180 students who appeared in the annual examination, 120 passed the math test, 90 passed the science test, and 60 passed both the math and science tests. How many did not pass either of the two tests?}
\lqn{104}{Out of the 180 students who appeared in the annual examination, 120 passed the math test, 90 passed the science test, and 60 passed both the math and science tests. How many passed the science test but not the math test?}
\lqn{105}{Out of the 180 students who appeared in the annual examination, 120 passed the math test, 90 passed the science test, and 60 passed both the math and science tests. How many failed the science test?}
\lqn{106}{Let $U=\{1,2,3,4,5,6,7,8,9,10\},\ A=\{2,4,6,8,10\},\ B=\{1,2,3,4,5\}$ and $C=\{1,3,5,7,9\}$ List the member of the following set: $A\cup C'$}
\lqn{107}{Let $U=\{1,2,3,4,5,6,7,8,9,10\},\ A=\{2,4,6,8,10\},\ B=\{1,2,3,4,5\}$ and $C=\{1,3,5,7,9\}$ List the member of the following set: $A'\cup C$}
\lqn{108}{A shopping mall has 100 products from various categories labeled 1 to 100, representing the universal st U. The products are categorized as follows: Set A: Electronics, consisting 30 products labeled from 1 to 30, Set B: Clothing comprises 25 products labeled from 31 to 55, Set C: Beauty products, comprising 25 products labeled from 76 to 100. Write each set in tabular form, and find the union of all three sets.}
\lqn{109}{Out of the 180 students who appeared in the annual examination, 120 passed the math test, 90 passed the science test, and 60 passed both the math and science tests. How many passed either the math or science test?}
\lqn{110}{Factorize: $2x^{2}+11x+15$}
\lqn{111}{Factorize: $3y^{2}-11y+6$}
\lqn{112}{Factorize: $2y^{2}-5y+2$}
\lqn{113}{Factorize: $4z^{2}-11z+6$}
\lqn{114}{Factorize: $6+7x-3x^{2}$}
\lqn{115}{Factorize: $2x^{2}+17x+26$}
\lqn{116}{Factorize the following expressions: $4x^{4}+81y^{4}$}
\lqn{117}{Factorize the following expressions: $a^{4}+64b^{4}$}
\lqn{118}{Factorize the following expressions: $x^{4}+4x^{2}+16$}
\lqn{119}{Factorize the following expressions: $x^{4}-14x^{2}+1$}
\lqn{120}{Factorize the following expressions: $x^{4}-30x^{2}y^{2}+9y^{4}$}
\lqn{121}{Factorize the following expressions: $x^{4}+11x^{2}y^{2}+y^{4}$}
\lqn{122}{Factorize the following expression: $(x+1)(x+2)(x+3)(x+4)+1$}
\lqn{123}{Factorize the following expression: $(x+2)(x-7)(x-4)(x-1)+17$}
\lqn{124}{Factorize the following expression: $(2x^{2}+7x+3)(2x^{2}+7x+5)+1$}
\lqn{125}{Factorize the following expression: $(3x^{2}+5x+3)(3x^{2}+5x+5)-3$}
\lqn{126}{Factorize the following expression: $(x+1)(x+2)(x+3)(x+6)-3x^{2}$}
\lqn{127}{Factorize the following expression: $(x+1)(x-1)(x+2)(x-2)+13x^{2}$}
\lqn{128}{Factorize: $a^{4}+64$}
\lqn{129}{Factorize: $(x^{2}+5x+4)(x^{2}+5x+6)-3$}
\lqn{130}{Factorize: $(x-2)(x+2)(x+1)(x-4)+2x^{2}$}
\lqn{131}{Factorize: $8x^{3}+60x^{2}+150x+125$}
\lqn{132}{Factorize: $x^{4}+x^{2}+25$}
\lqn{133}{Factorize: $x^{4}+y^{4}$}
\lqn{134}{Find HCF by factorization method: $a^{3}+2a^{2}-3a,\ 2a^{3}+5a^{2}-3a$}
\lqn{135}{Find HCF by factorization method: $t^{2}+3t-4,\ t^{2}+5t+4,\ t^{2}-1$}
\lqn{136}{Find HCF by factorization method: $x^{2}+15x+56,\ x^{2}+5x-24,\ x^{2}+8x$}
\lqn{137}{Find HCF of the following expressions by using division method: $27x^{3}+9x^{2}-3x-9,\ 3x-2$}
\lqn{138}{Find HCF of the following expressions by using division method: $x^{3}-9x^{2}+21x-15,\ x^{2}-4x+3$}
\lqn{139}{Find HCF of the following expressions by using division method: $2x^{3}+2x^{2}+2x+2,\ 6x^{3}+12x^{2}+6x+12$}
\lqn{140}{Find HCF of the following expressions by using division method: $2x^{3}-4x^{2}+6x,\ x^{3}-2x,\ 3x^{2}-6x$}
\lqn{141}{Find LCM of the following expressions by using prime factorization method: $x^{4}-16,\ x^{3}-4x$}
\lqn{142}{Find LCM of the following expressions by using prime factorization method: $16-4x^{2},\ x^{2}+x-6,\ 4-x^{2}$}
\lqn{143}{The HCF of two polynomials is y-7 and their LCM is $y^{3}-10y^{2}+11y+70$. If one of the polynomials is $y^{2}-5y-14$, find the other.}
\lqn{144}{The LCM and HCF to two polynomials p(x) and q(x) are $36x^{3}(x+a)(x^{3}-a^{3})$ and $x^{2}(x-a)$ respectively. If $p(x)=4x^{2}(x^{2}-a^{2})$, find q(x) and q(x).}
\lqn{145}{The HCF and LCM of two polynomials are (x+a) and $12x^{2}(x+a)(x^{2}-a^{2})$ respectively. Find the product of the two polynomials.}
\lqn{146}{Find the HCF by factorization method $x^{2}-27,\ x^{2}+6x-27,\ x^{2}-9$}
\lqn{147}{Find HCF of $6x^{3}-17x^{2}-5x+6$ and $6x^{3}-5x^{2}-3x+2$ by using division method.}
\lqn{148}{Find the LCM of the polynomials $x^{2}-3x+2,\ x^{2}-1$ and $x^{2}-5x+4$}
\lqn{149}{LCM and HCF of two polynomials are $x^{3}-10x^{2}+11x+70$ and x-7.}
\lqn{150}{The LCM of $x^{2}y+xy^{2}$ and $x^{2}$+xy is xy(x+y). Find the HCF}
\lqn{151}{Find the square root of the following polynomials by factorization method: $200t^{2}-120t+18$}
\lqn{152}{Find the square root of the following polynomials by factorization method: $40x^{2}+120x+90$}
\lqn{153}{Find the square root of the following polynomials by division method: $4x^{4}-28x^{3}+37x^{2}+42x+9$}
\lqn{154}{The potential energy V(x) in an electric field varies as a cubic function of distance x, given by: $V(x)=2x^{3}-6x^{2}+4x$. Determine where the potential energy is zero.}
\lqn{155}{In structural engineering, the deflection Y(x) of a beam is given by: $Y(x)=2x^{2}-8x+6$. This equation gives the vertical deflection at any point x along the beam. Find the points of zero deflection.}
\lqn{156}{Find the square root of the following polynomials by division method: $121x^{4}-198x^{3}-183x^{2}+216x+144$}
\lqn{157}{Find the square root of the following polynomials by division method: $x^{4}-10x^{3}y+27x^{2}y^{2}-10xy^{3}+y^{4}$}
\lqn{158}{Find the square root of the following polynomials by division method: $x^{4}-10x^{3}y+27x^{2}y^{2}-10xy^{3}+y^{4}$}
\lqn{159}{An investor's return R(x) in rupees after investing x thousand rupees is given by the quadratic expression: $R(x)=-x^{2}+6x-8$. Factorize the expression and find the investment levels that result in zero return.}
\lqn{160}{A company's profit P(x) in rupees from selling x units of a products is modeled by the cubic expression: $P(x)=x^{3}-15x^{2}+75x-125$. Find the break-even point(s), where the profit is zero.}
\lqn{161}{Find the square root of the expression $36x^{4}-36x^{2}+9$}
\lqn{162}{Find the square root of the polynomial $x^{4}-12x^{3}+42x^{2}-36x+9$}
\lqn{163}{Cost function for producing a part is modeled by: $C(x)=5x^{2}-25x+30$. Where x is the width of the component and C(x) is the cost. Find the value of x where C(x) is minimum.}
\lqn{164}{A company's profit P(x) is modeled by the quadratic equation: $P(x)=-5x^{2}+50x-120$. Where x represent the number of units produced and P(x) represent the profit in dollars. Find how many units should be produced to maximize profit.}
\lqn{165}{The potential energy U(x) of a particle moving in a cubic potential is expressed as: $U(x)=x^{3}-6x^{2}+12x-8$. Factorize the expression to find the points where the energy in minimized.}
\lqn{166}{Factorize the following: (x+3)(x+4)(x+5)(x+6)-360}
\lqn{167}{Factorize the following: $(x^{2}+6x+3)(x^{2}+6x-9)+36$}
\lqn{168}{Find LCM and HCF by prime factorization method: $4x^{3}+12x^{2},\ 8x^{2}+16x$}
\lqn{169}{Find LCM and HCF by prime factorization method: $x^{3}+3x^{2}-4x,\ x^{2}-x-6$}
\lqn{170}{Huria is analyzing the total cost of her loan, modeled by the expression $C(x)=x^{2}-8x+15$, where x represent the number of years. what is the optimal repayment period for Huria's loan?}
\lqn{171}{Solve and represent the solution on a real line: $\dfrac{2x-1}{3}-\dfrac{3x}{4}=\dfrac{5}{6}$}
\lqn{172}{Solve and represent the solution on a real line: $\dfrac{-5x}{10}=9-\dfrac{10}{5}x$}
\lqn{173}{Solve following inequality and represent the solution on a real line: $\dfrac{5}{3}x-\dfrac{3}{4}<\dfrac{-1}{12}$}
\lqn{174}{Solve following inequality and represent the solution on a real line: $\dfrac{1}{4}x-\dfrac{1}{2}\leq-1+\dfrac{1}{2}x$}
\lqn{175}{Indicate the solution region of the following linear inequalities by shading: $2x-3y\leq6,\ 2x-3y\leq12$}
\lqn{176}{Indicate the solution region of the following linear inequalities by shading: $x+y\geq5,\ -y+x\leq1$}
\lqn{177}{Indicate the solution region of the following linear inequalities by shading: $3x+7y\geq21,\ x-y\leq2$}
\lqn{178}{Indicate the solution region of the following linear inequalities by shading: $4x-3y\leq12,\ x\geq-\dfrac{3}{2}$}
\lqn{179}{Indicate the solution region of the following linear inequalities by shading: $3x+7y\geq21,\ y\leq4$}
\lqn{180}{Indicate the solution region of the following linear inequalities by shading: $5x+7y\leq35,\ x-2y\leq2$}
\lqn{181}{Find the solution region by drawing the graph of the system of inequalities: $x-2y\leq6,\ 2x+y\geq2$}
\lqn{182}{Find minimum and maximum values of z = 3x + y; subject to the constraints: $3x+5y\geq15,\ x+6y\geq9,\ x\geq0,\ y\geq0$}
\lqn{183}{Maximize f(x,y)=2x+5y, subject to the constraints: $2y-x\leq8,\ x-y\leq4,\ x\geq0,\ y\geq0$}
\lqn{184}{Maximize f(x,y)=x+3y ; subject to constraints $2x+5y\leq30,\ 5x+4y\leq20,\ x\geq0,\ y\geq0$}
\lqn{185}{Maximize z = 2x + 3y; subject to constraints: $2x+y\leq4,\ 4x-y\leq4,\ x\geq0,\ y\geq0$}
\lqn{186}{Minimize z = 2x + y; subject to the constraints: $x+y\geq3,\ 7x+5y\leq35,\ x\geq0,\ y\geq0$}
\lqn{187}{Maximize the function defined as: f(x,y)=2x+3y subject to the constraints: $2x+y\leq8,\ x+2y\leq14,\ x\geq0,\ y\geq0$}
\lqn{188}{Shade the feasible region and find the corner points for the following system of inequalities (or subject to the following constraints): $x-y\leq3,\ x+2y\leq6,\ x\geq0,\ y\geq0$}
\lqn{189}{Find the maximum and minimum values of the function defined as: $f(x,y)=2x+3y$ subject to the constrains; $x-y\leq2,\ x+y\leq4,\ x\geq0,\ y\geq0$.}
\lqn{190}{Solve and represent their solutions on real line: $5(x-3)\geq26x-(10x+4)$.}
\lqn{191}{Find the solution region of the following linear equalities: $3x-4y\leq12,\ 3x+2y\geq3$}
\lqn{192}{Find the solution region of the following linear equalities: $2x+y\leq4,\ x+2y\leq6$}
\lqn{193}{Find the maximum value of g(x,y)= x+4y subject to constraints $x+y\leq4,\ x\geq0$ and $y\geq0$.}
\lqn{194}{Find the minimum value of f(x,y) =3x+5y subject to constrains $x+3y\geq3,\ x+y\geq2,\ x\geq0,\ y\geq0$}
\lqn{195}{If the central angle of a sector is 60$^{o}$ and the radius of the circle is 12 cm, find the area of the sectors and the percentage of the total area of the circle it represents.}
\lqn{196}{A circular sector of radius r = 12 cm has an angle of 150$^{o}$. This sector is cut out and then bent to form a cone. What is the slant height and the radius of the base pf this cone.? Hint Arc length of sector = circumference of cone.}
\lqn{197}{Prove the following trigonometric identities: $\dfrac{\sin\theta}{\mathrm{cosec}\,\theta}+\dfrac{\cos\theta}{\sec\theta}=1$}
\lqn{198}{Prove the following trigonometric identities: $\dfrac{1-\sin\theta}{\cos\theta}=\dfrac{\cos\theta}{1+\sin\theta}$}
\lqn{199}{Prove the following trigonometric identities: $(\sec\theta-\tan\theta)^{2}=\dfrac{1-\sin\theta}{1+\sin\theta}$}
\lqn{200}{Prove the following trigonometric identities: $(\tan\theta+\cot\theta)^{2}=\sec^{2}\theta\,\mathrm{cosec}^{2}\theta$}
\lqn{201}{Prove the following trigonometric identities: $\dfrac{\tan\theta+\sec\theta-1}{\tan\theta-\sec\theta+1}=\tan\theta+\sec\theta$}
\lqn{202}{Prove the following trigonometric identities: $\sin^{3}\theta-\cos^{3}\theta=(\sin\theta-\cos\theta)(1+\sin\theta\cos\theta)$}
\lqn{203}{Prove the following trigonometric identities: $\sin^{6}\theta-\cos^{6}\theta=(\sin^{2}\theta-\cos^{2}\theta)(1-\sin^{2}\theta\cos^{2}\theta)$}
\lqn{204}{Show that: $\dfrac{1}{\mathrm{cosec}\,\theta-\cot\theta}-\dfrac{1}{\sin\theta}=\dfrac{1}{\sin\theta}-\dfrac{1}{\mathrm{cosec}\,\theta+\cot\theta}$}
\lqn{205}{Show that: $\sin^{6}\theta+\cos^{6}\theta=1-3\sin^{2}\theta\cos^{2}\theta$}
\lqn{206}{If $\tan\theta=\dfrac{3}{4}$, find the remaining trigonometric ratios, when $\theta$ lies in first quadrant.}
\lqn{207}{Evaluate: $\sin60^{o}\cos30^{o}+\cos60^{o}\sin30^{o}$}
\lqn{208}{If $\sin\dfrac{\pi}{4}$ and $\cos\dfrac{\pi}{4}$ equal to $\dfrac{1}{\sqrt{2}}$ each, then find the value of the following: $2\sin45^{o}-2\cos45^{o}$}
\lqn{209}{If $\sin\dfrac{\pi}{4}$ and $\cos\dfrac{\pi}{4}$ equal to $\dfrac{1}{\sqrt{2}}$ each, then find the value of the following: $3\cos45^{o}+4\sin45^{o}$}
\lqn{210}{If $\sin\dfrac{\pi}{4}$ and $\cos\dfrac{\pi}{4}$ equal to $\dfrac{1}{\sqrt{2}}$ each, then find the value of the following: $5\cos45^{o}-3\sin45^{o}$}
\lqn{211}{Solve the following triangle when $m\angle B=90^{\circ}$, $m\angle C=60^{\circ}$, $c=3\sqrt{3}\ cm$}
\lqn{212}{Solve the following triangle when $m\angle B=90^{\circ}$: a = 12 cm, C= 6 cm}
\lqn{213}{Solve the following triangle when $m\angle B=90^{\circ}$: b = 10cm, a = 6 cm}
\lqn{214}{Calculate the length x in the adjoining figure.}
\pfig{fig214}{5.5cm}
\lqn{215}{Solve the following triangle when $m\angle B=90^{\circ}$, $m\angle C=45^{\circ}$, a = 8 cm}
\lqn{216}{Solve the following triangle when $m\angle B=90^{\circ}$, $m\angle A=60^{\circ}$, c = 4 cm}
\lqn{217}{Solve the following triangle when $m\angle B=90^{\circ}$, $m\angle A=60^{\circ}$, c = 4 cm}
\lqn{218}{Let Q and R be the two points on the same bank of a canal. The point P is placed on the other bank straight to point R. Find the width of the canal and the angle PQR. image}
\lqn{219}{If the ladder is placed along the wall such that the foot of the ladder is 2 m away from the wall. If the length of the ladder is 8 m, find the height of the wall.}
\lqn{220}{The diagonal of a rectangular field ABCD is (x + 9)m and the sides are (x + 7)m and x m. Find the value of x.}
\pfig{fig220}{4.5cm}
\lqn{221}{Solve triangle ABC, in which $m\angle B=90^{\circ}$, $m\angle A=30^{\circ}$, a = 2.}
\lqn{222}{Solve triangle ABC, when $a=\sqrt{2}\ cm$, c = 1 cm $m\angle B=90^{\circ}$}
\lqn{223}{An isosceles triangle has a vertical angle of 120$^{o}$ and a base 10 cm long. Find the length of its altitude.}
\lqn{224}{A light house tower is 150 m high from the sea level. The angle of depression from the top of the tower to a ship is 60$^{o}$. Find the distance between the ship and the tower.}
\lqn{225}{Measure of an angle of elevation of the top of a pole is 15$^{o}$ from a point on the ground, in walking 100 m towards the pole the measure of angle is found to be 30$^{o}$. Find the height of the pole.}
\lqn{226}{Measure of angle of elevation of the top of a cliff is 25$^{o}$. on walking 100 metres towards the cliff, measure the angle of elevation of the top is 45$^{o}$. Find the height of the cliff.}
\lqn{227}{From the top of a hill 300 m high, the measure of the angle of depression of a point on the nearer shore of the river is 70$^{o}$ and measure of the angle of depression of a point, directly across the river is 50$^{o}$. Find the width of the river. How far is the river from the foot of the hill?}
\lqn{228}{A kite has 120 m of string attached to it when at an angle of elevation of 50$^{o}$. How far is it above the hand holding it? (Assume that the string is stretched.)}
\lqn{229}{From the top of a lookout tower, the angle of depression of a building has on the ground level of 45$^{o}$. How far is a man on the ground from the tower, if the height of the tower is 30 m?}
\lqn{230}{The angle of elevation of the top of a pole 40 m high is 60$^{o}$ when seen from a point on the ground level. Find the distance of the point from the foot of the pole.}
\lqn{231}{Prove the following identity: $\sin\theta\,(\mathrm{cosec}\,\theta-\sin\theta)=\dfrac{1}{\sec^{2}\theta}$}
\lqn{232}{Prove the following identity: $\dfrac{\mathrm{cosec}\,\theta-\sec\theta}{\mathrm{cosec}\,\theta+\sec\theta}=\dfrac{1-\cos\theta}{1+\tan\theta}$}
\lqn{233}{Prove the following identity: $\dfrac{\cos\theta+\sin\theta}{\cos\theta-\sin\theta}+\dfrac{\cos\theta-\sin\theta}{\cos\theta+\sin\theta}=\dfrac{2}{1-2\sin^{2}\theta}$}
\lqn{234}{Prove the following identity: $\dfrac{1+\cos\theta}{1-\cos\theta}=(\mathrm{cosec}\,\theta+\cot\theta)^{2}$}
\lqn{235}{If $\tan\theta=\dfrac{3}{\sqrt{2}}$ then find the remaining trigonometric ratio when $\theta$ lies in first quadrant.}
\lqn{236}{Show that: the points $A(0,2),\ B(\sqrt{3},1)$ and $C(0,-2)$ are vertices of a right triangle.}
\lqn{237}{Show that: the point $A(3,1),\ B(-2,-3)$ and $C(2,2)$ are vertices of an isosceles triangle.}
\lqn{238}{Show that: the points $A(5,2),\ B(-2,3),\ C(-3,-4)$ and $D(4,-5)$ are vertices of a parallelogram.}
\lqn{239}{Find h such that the points $A(\sqrt{3},-1),\ B(0,2)$ and $C(h,-2)$ are vertices of a right triangle with right angle at the vertex A.}
\lqn{240}{Find h such that $A(-1,h),\ B(3,2)$ and $C(7,3)$ are colinear.}
\lqn{241}{The points $A(-5,-2)$ and $B(5,-4)$ are ends of a diameter of a circle. Find the centre and radius of the circle.}
\lqn{242}{Find h such that the points $A(h,1),\ B(2,7)$ and $C(-6,-7)$ are vertices of a right triangle with right angle at the vertex A.}
\lqn{243}{A quadrilateral has the points $A(9,3),\ B(-7,7),\ C(-3,-7)$ and $D(5,-5)$ as its vertices. Find the midpoints of its sides. Show that the figure formed by joining the midpoints consecutively is a parallelogram.}
\lqn{244}{Find the distance between the points: $C(-4,-2),\ D(0,9)$}
\lqn{245}{Using slops, show that the triangle with its vertices $A(6,1),\ B(2,7)$ and $C(-6,-7)$ is a right triangle.}
\lqn{246}{Two pairs of points are given. Find whether the two lines determined by these points are: (i) parallel (ii) Perpendicular (iii) None. $(1,-2),\ (2,4)$ and $(4,1),\ (-8,2)$}
\lqn{247}{Two pairs of points are given. Find whether the two lines determined by these points are: (i) parallel (ii) Perpendicular (iii) None. $(-3,4),\ (6,2)$ and $(4,5),\ (-2,-7)$}
\lqn{248}{Find an equation of the perpendicular bisector of the segment joining the points $A(3,5)$ and $B(9,8)$.}
\lqn{249}{Find an equation of the line through $(-4,-6)$ and perpendicular to a line having slope $\dfrac{-3}{2}$}
\lqn{250}{Find an equation of the line through $(11,-5)$ and parallel to a line with slope -24.}
\lqn{251}{Convert the following equation into slope intercept from, two intercept form and normal form. $2x-4y+11=0$}
\lqn{252}{Convert the following equation into slope intercept from, two intercept form and normal form. $4x+7y-2=0$}
\lqn{253}{Convert the following equation into slope intercept from, two intercept form and normal form. $15y-8x+3=0$}
\lqn{254}{In the following check whether the two lines are: (i) parallel (ii) perpendicular (iii) neither parallel nor perpendicular $2x+y-3=0$; $4x+2y+5=0$}
\lqn{255}{In the following check whether the two lines are: (i) parallel (ii) perpendicular (iii) neither parallel nor perpendicular $3y=2x+5$; $3x+2y-8=0$}
\lqn{256}{In the following check whether the two lines are: (i) parallel (ii) perpendicular (iii) neither parallel nor perpendicular $4y+2x-1=0$; $x-2y-7=0$}
\lqn{257}{Find an equation of the line through $(-4,7)$ and parallel to the line $2x-7y+4=0$.}
\lqn{258}{Find an equation of the line through $(5,-8)$ and perpendicular to the join of $A(-15,-8),\ B(10,7)$.}
\lqn{259}{Find the equation of the straight line if: its slope is 2 and y-intercept is 5.}
\lqn{260}{Find the equation of the straight line if: it is perpendicular to a line with slope -6 and its y-intercept is $\dfrac{4}{3}$}
\lqn{261}{Write down an equation of the straight line passing through (5,1) and parallel to a line passing through the points: (0,1), (7,15).}
\lqn{262}{The length of perpendicular from the origin to a line is 5 units and the inclination of this perpendicular is 120$^{o}$. Find the slope and y-intercept of the line.}
\lqn{263}{An architect is designing a park with two building located at $(10,8)$ and $(4,3)$ on the grid. Calculate the straight-line distance between the buildings. Assume the coordinates are in meters.}
\lqn{264}{A delivery driver needs to calculate the distance between two delivery locations. One location is at $(7,2)$ and the other at $(12,10)$ on the city grid map, where each unit represent kilometres. What is the distance between the two locations?}
\lqn{265}{A ship is navigation from port A located at $(12^{\circ}N,\ 65^{\circ}W)$ to port B at $(20^{\circ}N,\ 45^{\circ}W)$. If the ship travels along the shortest path on the surface of the Earth, calculate the straight line distance between the points.}
\lqn{266}{Farah is fencing around a rectangular field with corners at $(0,0),\ (0,5),\ (8,5)$ and $(8,0)$. How much fencing material will she need to cover the entire perimeter of the field?}
\lqn{267}{An airplane is flying from city X at $(40^{\circ}N,\ 100^{\circ}W)$ to city Y at $(50^{\circ}N,\ 80^{\circ}W)$. Use coordinate geometry, calculate the shortest distance between these two cities.}
\lqn{268}{A land surveyor is marking out a rectangular plot of land with corners at $(3,1),\ (3,6),\ (8,6)$, and $(8,1)$. Calculate the perimeter.}
\lqn{269}{A landscaper needs to install a fence around a rectangular garden. The garden has its corners at the coordinates: $A(0,0),\ B(5,0),\ C(5,3)$ and $D(0,3)$. How much fencing is required?}
\lqn{270}{A landscaper is designing a triangular garden with corners at points A(2,3), B(5,7) and C(6,2). Calculate the lengths of the sides of the triangle.}
\lqn{271}{Find the question of the line in the form y = mx + c that passes through the points (3,7) and (5,11).}
\lqn{272}{A drone is flying from point (2,3) to point (10,15) on the grid. Calculate the gradient of the line along which the drone is flying and the total distance travelled.}
\lqn{273}{Formulate simple deductive proof's of the following algebraic expression, prove that L.H.S is equal to the R.H.S. $(x-4)^{2}+9=x^{2}-8x+25$}
\lqn{274}{Formulate simple deductive proof's of the following algebraic expression, prove that L.H.S is equal to the R.H.S. $(x+1)^{2}-(x-1)^{2}=4x$}
\lqn{275}{Formulate simple deductive proof's of the following algebraic expression, prove that L.H.S is equal to the R.H.S. $(x+5)^{2}-(x-5)^{2}=20x$}
\lqn{276}{Prove the following by justifying each step: $\dfrac{4+16x}{4}=1+4x$}
\lqn{277}{Prove the following by justifying each step: $\dfrac{6x^{2}+18x}{3x^{2}-9}=\dfrac{2x}{x-3}$}
\lqn{278}{Prove the following by justifying each step: $\dfrac{x^{2}+7x+10}{x^{2}-3x-9}=\dfrac{x+5}{x-5}$}
\lqn{279}{Prove the following statement: If x is an odd integer, then show that $x^{2}-4x+6$ is odd.}
\lqn{280}{Prove the following statement: If x is an even integer then show that $x^{2}+2x+4$ is even.}
\lqn{281}{Prove that for any two non-empty sets A and B, $(A\cap B)'=A'\cup B'$}
\lqn{282}{Construct the truth table of: $[(p\rightarrow q)\wedge p]$ and $[(p\rightarrow q)\wedge p]\rightarrow q$}
\lqn{283}{Prove that for any two non-empty sets A and B, $(A\cup B)'=A'\cap B'$.}
\lqn{284}{Prove that $\dfrac{a}{b}+\dfrac{c}{d}=\dfrac{ad+bc}{bd}$ where a,b,c and d are non-zero real numbers.}
\lqn{285}{Prove that: $(x+1)^{2}+7=x^{2}+2x+8$}
\lqn{286}{Prove that $\dfrac{45x+15}{15}=3x+1$ by justifying each step.}
\lqn{287}{A plank is placed straight upstairs that 20 cm wide and 16 cm deep. A rectangular box of height 8 cm and width x cm is placed on a stair under the plank. Find the value of x: (image)}
\lqn{288}{A man who is 1.8 m tall casts a shadow of a 0.76 m in length. If at the same time a telephone pole costs a 3 m shadow, Find the height of the pole.}
\lqn{289}{Draw an isosceles trapezoid ABCD where $\overline{AB}\parallel\overline{CD}$ and $m\overline{AB}>m\overline{CD}$. Draw diagonals $\overline{AC}$ and $\overline{BD}$, intersecting at E. Prove that $\triangle ABE$ is similar to $\triangle CDE$. If $m\overline{AB}=8\ cm$, $m\overline{CD}=4\ cm$ and $m\overline{AE}=3\ cm$ find the length of $\overline{CE}$.}
\lqn{290}{A regular dodecagon has its side lengths decreased by a factor of $\dfrac{1}{\sqrt{2}}$. If the perimeter of the original dodecagon is 72 cm. What is the side length of scaled dodecagon?}
\lqn{291}{In the given figure, $\triangle ABC\sim\triangle DEF$, $m\overline{AB}=12cm,\ m\overline{AC}=20cm,\ m\overline{BC}=16cm$. In $\triangle DEF,\ m\overline{DE}=6cm$ Find $m\overline{DE}$ and $m\overline{EF}$}
\pfig{fig291}{5.5cm}
\lqn{292}{Find whether the parallelograms ara similar given that one of the angle between sides is 45$^{o}$ in both parallelograms.}
\lqn{293}{In the triangles XBC and XDE, find the value of x and y.}
\lqn{294}{The quadrilateral ABCD has side lengths $m\overline{AB}=5cm,\ m\overline{BC}=8,\ m\overline{CD}=10\ cm,\ m\overline{AD}=12\ cm$, and angles are $m\angle A=90^{\circ},\ m\angle B=120^{\circ}$ and $m\angle C=90^{\circ}$. Quadrilateral EFGH has side lengths $m\overline{EF}=10cm,\ m\overline{FG}=16cm,\ m\overline{GH}=20cm,\ m\overline{EH}=24cm$ and angles are $m\angle E=90^{\circ},\ m\angle F=120^{\circ}$ and $m\angle H=60^{\circ}$. Prove that the quadrilateral ABCD is similar to the quadrilateral EFGH.}
\lqn{295}{Given that area of $\triangle ABC=36\ cm^{2}$ and $m\overline{AB}=6\ cm$, $m\overline{BD}=4\ cm$. Find the area of $\triangle ADE$}
\lqn{296}{Given that area of $\triangle ABC=36\ cm^{2}$ and $m\overline{AB}=6\ cm$, $m\overline{BD}=4\ cm$. the area of trapezium BCED}
\lqn{297}{Given that $\triangle ABC$ and $\triangle DEF$ are similar, with a scale factor of k=3. If the area of $\triangle ABC$ is 50 $cm^{2}$, find the area of triangle $\triangle DEF$.}
\lqn{298}{Quadrilaterals ABCD and EFGH are similar, with a scale factor of $k=\dfrac{1}{4}$. If the area of quadrilaterals ABCD is 64$cm^{2}$, find the area of quadrilaterals EFGH.}
\lqn{299}{The areas of two similar triangles are 16 $cm^{2}$ and 25 $cm^{2}$. What is the ratio of a pair of corresponding sides?}
\lqn{300}{The areas of two similar triangles are 144 $cm^{2}$ and 81 $cm^{2}$. If the base of the large triangle is 30 cm, find the corresponding base of the smaller triangle.}
\lqn{301}{A regular heptagon is inscribed in a large regular heptagon and each side of the larger heptagon is 1.7 times the side of the smaller heptagon. If the area of the smaller heptagon is known to be 100 $cm^{2}$, find the area of the larger heptagon.}
\lqn{302}{The ratio of the heights of two similar cylindrical water tanks is 5:3: 1. if the surface area of the larger tank is 250 square metres, find the surface area of the smaller tank. 2. If the volume of the smaller tank is 270 cubic metres, find the volume of the larger tank.}
\lqn{303}{The ratio of the corresponding lengths of two similar canonical cans is 3:2. (i) The larger canonical can have surface area of 96 m$^{2}$. Find the surface area of the smaller canonical can. (ii) The smaller canonical can have a volume of 240 m$^{3}$. Find the volume of larger canonical can.}
\lqn{304}{The ratio of the corresponding lengths of two similar cylindrical cans is 3:2: The larger cylindrical can has surface area of 67.5 square metres. Find the surface area of the smaller cylindrical can.}
\lqn{305}{The ratio of the corresponding lengths of two similar cylindrical cans is 3:2: The smaller cylindrical can has a volume of 132 cubic metres. Find the volume of larger tin can.}
\lqn{306}{A solid cone C is cut into two pieces A and B with sloping edges 6 cm and 4 cm. find the ratio of: the diameter of the base of the cones A and C. (image)}
\lqn{307}{A solid cone C is cut into two pieces A and B with sloping edges 6 cm and 4 cm. find the ratio of: The volumes of the two cones A and C. (image)}
\lqn{308}{A solid cone C is cut into two pieces A and B with sloping edges 6 cm and 4 cm. find the ratio of: the area of the base of the cones A and C. (image)}
\lqn{309}{A solid cone C is cut into two pieces A and B with sloping edges 6 cm and 4 cm. find the ratio of: If volume of cone A is 72 cm$^{3}$, find the volume of solid B. (image)}
\lqn{310}{The mass of a sack of rice is 50 kg and height 4 m. Find the mass of the similar sack of rice with height of 6m.}
\lqn{311}{A shape is created by a cutting a square in half diagonally and then attaching a right-angled triangle to the hypotenuse of each half. Explain why this shape can tessellate and calculate the interior angle of the new shape.}
\lqn{312}{A tessellation is created by repeatedly reflecting a basic shape. The basic shape is a right-angled triangle with sides of lengths 3, 4 and 5 units. Find: The minimum number of reflections needed to create a tessellation that covers a square with an area of 3600 square units.}
\lqn{313}{A tessellation is created using regular hexagons. Each hexagons has a side length of 5 cm. Find the total area of the tessellation if it consists of 25 hexagons and total perimeter of the outer edge of the tessellation, assuming it's perfect hexagon.}
\lqn{314}{In a parallelogram ABCD if $m\angle DAB=70^{\circ}$, find the measures of all other angles in the parallelogram.}
\lqn{315}{A window has a trapezoidal shape with parallel sides of 3 m and 1.5 m and a height of 2 m. Find the area of the window.}
\lqn{316}{Falak plans to tile a square patio with an area of 100 square metres. He decides to use both square tiles and triangular tiles, each with an area of 0.25 square metres, if 60\% of the tiles will be square and 40\% will be trangular, how many tiles of each shape are needed?}
\lqn{317}{Ali wants to create a floor design that uses regular hexagons (each with a side length of 1 metre) and equilateral triangle (each with a side length of 1 metre) to cover a rectangular area measuring 10 m by 5 m. Find how many hexagons and triangles Ali will need to complete the tessellation.}
\lqn{318}{A parallelogram-shaped room has a base of 10 meters and a height of 8 m. Babar wants to carpet the room using rolls that cover 20 m$^{2}$ each. How many rolls of carpet do he need?}
\lqn{319}{A tessellation is created using a combination of regular pentagons and decagons. Find the sum of the angles at a vertex where a pentagon and a decagon meet.}
\lqn{320}{Plot the graph of the following quadratic and cubic function: $y=x^{3}+2x^{2}-5x-6;\ -3.5\leq x\leq2.5$}
\lqn{321}{Plot the graph of the following quadratic and cubic function: $y=x^{2}+x-2$}
\lqn{322}{Plot the graph of the following quadratic and cubic function: $y=x^{3}+3x^{2}+2x;\ -2.5\leq x\leq0.5$}
\lqn{323}{Plot the graph of the following quadratic and cubic function: $y=5x^{2}-2x-3$}
\lqn{324}{Plot the graph of $y=2x^{2}-4x+3$ from -1 to 3. Draw the tangent at (2,3) and find the gradient.}
\lqn{325}{Plot the graph of $y=3x^{2}+x+1$ and draw tangent at (1,5). Also find gradient of the tangent line at this point.}
\lqn{326}{The strength of students in a school is 1000 in 2016. If the strength decay according to the equation S= 1000 $e^{-1}$, where S is the number of student at time t. Graph the given equation t=0 (in 2016) to t=9 (in 2025).}
\lqn{327}{The strength of students in a school is 1000 in 2016. If the strength decay according to the equation S= 1000 $e^{-1}$, where S is the number of student at time t. from the graph, estimate the student's strength in 2019 and 2023.}
\lqn{328}{A newspaper agency fixed cost of Rs. 70 per edition and marginal printing and distribution costs of Rs. 40 per copy. Profit function is $p(x)=0.10x+70$, where x is the number of newspapers. Plot the graph and find the profit for 500 newspapers.}
\lqn{329}{Ali manufactures expensive shirts for sale to a school. Its cost (in rupees) for x shirts is $C(x)=1500+10x+0.2x^{2},\ 0\leq x\leq150$. Plot the graph and find the cost of 200 shirts.}
\lqn{330}{The population of a village was 753 in 2010. If the population grows according to the equation $p=753e^{0.03t}$, where p is the number of person in the population at time t. Graph the population equation for t=0 (in 2010) to t = 30 (in 2040).}
\lqn{331}{The population of a village was 753 in 2010. If the population grows according to the equation $p=753e^{0.03}$, where p is the number of person in the population at time t. from the graph, estimate the population (i) in 2020 and (ii) in 2030.}
\lqn{332}{Sketch the graph of $y=x^{2}-3x-2$ for values of x from -8 to 8, draw a tangent line at x = -6 and determine the gradient.}
\lqn{333}{Majid's salary S(x) in rupees is based on the following formula: S(x)= 25000 + 1500x, where x is the number of years he worked. Sketch and interpret the graph of salary function for $0\leq x\leq10$.}
\lqn{334}{A company manufactures footballs. The cost of manufacturing x footballs is C(x)= 90,000 + 600x. The revenue from the selling is R(x)= 1,800x. Find the break-even point and determine the profit or loss when 200 footballs are sold. Draw the graph of both the functions and identify the break-point.}
\lqn{335}{Sales of a new magazine are expected to grow according to the equation: S=200000 $\left(1-e^{-0.05t}\right)$, where t is given in weeks. Plot graph of sales for the first 50 weeks.}
\lqn{336}{Sales of a new magazine are expected to grow according to the equation: S=200000 $\left(1-e^{-0.05t}\right)$, where t is given in weeks. Calculate the number of magazines sold, when t=5 and t=35.}
\lqn{337}{Plot the graph of $y=\dfrac{1}{2}(x+4)(x-1)(x-3)$ from -5 to 4.}
\lqn{338}{The supply and demand functions for a particular market are given by the equation: $P_{s}=Q^{2}+5$ and $P_{d}=Q^{2}-10Q$, where P represents price and Q represent quantity, Sketch the graph if each function over the interval Q=-20 to Q=20.}
\lqn{339}{A television manufacturer company make 40 inches LEDS. The cost of manufacturing x LEDs is C(x)= 60,000 + 250 x and the revenue from selling x LEDs is R(x)=1200x. Find the break-even point and find the profit or loss when 100 LEDs are sold. Identify the break-even point graphically.}
\lqn{340}{Verify that the angle bisectors of $\triangle ABC$ are concurrent with the following measurement: $m\overline{AB}=4.5\ cm,\ m\angle A=45^{\circ}$ and $m\overline{AC}=5.3\ cm$.}
\lqn{341}{Verify that the angle bisectors of $\triangle ABC$ are concurrent with the following measurement: $m\overline{AB}=6cm,\ m\angle A=150^{\circ}$ and $m\angle B=60^{\circ}$.}
\lqn{342}{Construct $\triangle ABC$ with the given measurements and verify that the perpendicular bisectors of the triangle are concurrent: $m\overline{AB}=5\ cm,\ m\overline{BC}=6cm$ and $m\overline{AC}=7\ cm$}
\lqn{343}{Construct $\triangle ABC$ with the given measurements and verify that the perpendicular bisectors of the triangle are concurrent: $m\overline{AB}=7.1\ cm,\ m\angle B=135^{\circ}$ and $m\overline{BC}=6.5\ cm$.}
\lqn{344}{Construct $\triangle LMN$ of the following measurements and verify that the medians of the triangle are concurrent: $m\overline{LM}=4.9\ cm,\ m\angle L=51^{\circ}$ and $m\angle M=38^{\circ}$}
\lqn{345}{Construct $\triangle LMN$ of the following measurements and verify that the medians of the triangle are concurrent: $m\overline{MN}=4.8\ cm,\ m\angle N=30^{\circ}$ and $m\overline{LM}=8.1cm$.}
\lqn{346}{Verify that the angle bisectors of $\triangle ABC$ are concurrent with the following measurement: $m\overline{AB}=4.5\ cm,\ m\angle A=45^{\circ}$ and $m\overline{AC}=5.3\ cm$.}
\lqn{347}{Given the measurement of $\triangle DEF$: $m\overline{DE}=4.8\ cm,\ m\overline{EF}=4\ cm$ and $m\angle E=45^{\circ}$, draw altitudes of $\triangle DEF$ and find orthocentre.}
\lqn{348}{Construct the following triangles and find whether there exists any ambiguous case: $\triangle BCD$ ; $m\overline{BC}=5\ cm,\ m\angle B=62^{\circ}$ and $m\overline{CD}=4.7\ cm$.}
\lqn{349}{Construct the following triangles and find whether there exists any ambiguous case: $\triangle KLM$ ; $m\overline{LM}=6\ cm,\ m\angle M=42^{\circ}$ and $m\overline{LN}=5\ cm$.}
\lqn{350}{Construct a triangle of sides 5.3 cm, 5.9 cm and 6.2 cm.}
\lqn{351}{Construct a triangle BCD in which measures of two sides are 5.5 cm and 4.2 cm and measure of their included angle is 60$^{0}$.}
\lqn{352}{Draw a triangle CDE when $m\overline{DE}=4.3\ cm,\ m\angle D=30^{\circ}$ and $m\angle E=120^{\circ}$.}
\lqn{353}{Construct a triangle DEF when $m\overline{DE}=6\ cm,\ m\angle D=110^{\circ}$ and $m\overline{EF}=9\ cm$.}
\lqn{354}{Construct triangles DEF and $DEF'$ when $m\overline{DE}=6\ cm,\ m\angle D=30^{\circ}$ and $m\overline{EF}=3.6\ cm$.}
\lqn{355}{In the above example if we take: $m\overline{EF}=3\ cm$}
\lqn{356}{In the above example if we take: $m\overline{EF}=2.5\ cm$}
\lqn{357}{Draw perpendicular bisector of the triangle EFG with $m\overline{EF}=5\ cm,\ m\overline{FG}=2.5\ cm$ and $m\overline{EG}=4.3\ cm$.}
\lqn{358}{Draw angle bisector of a triangle FGH if: $m\overline{FG}=5.2,\ m\overline{GH}=4.1$ and $m\angle FGH=120^{\circ}$.}
\lqn{359}{Construct triangle GHI in which $m\overline{GH}=5.7,\ m\angle G=68^{\circ}$ and $m\angle G=50^{\circ}$. Prove that altitudes of the $\triangle GHI$ are concurrent.}
\lqn{360}{Construct a CDE with $m\overline{CD}=7.6\ cm,\ m\angle D=45^{\circ}$ and $m\overline{DE}=5.9\ cm$. Draw the locus of all points which are: equidistant from C and D. Mark the point X where the two loci intersect.}
\lqn{361}{Two point A and B are 8.2 cm apart. Construct the locus of points 5 cm from point A.}
\lqn{362}{Construct a locus of point 2.2 cm from line segment CD of measure 5.7 cm.}
\lqn{363}{Construct an angle ABC=105$^{o}$. Construct a locus of a point P which shows moves such that it is equidistant from $\overline{BA}$ and $\overline{BC}$.}
\lqn{364}{Two point E and F are 5.4 cm apart. Construct a locus of a point P which moves such that it is equidistant E and F.}
\lqn{365}{The island has two main cities A and B 8 km apart. Kashif lives on the island exactly 6.8 km from city A and exactly 7.3 km from city B. Mark with a cross the point on the island where Kashif could line.}
\lqn{366}{Construct a CDE with $m\overline{CD}=7.6\ cm,\ m\angle D=45^{\circ}$ and $m\overline{DE}=5.9\ cm$. Draw the locus of all points which are: equidistant from $\overline{CD}$ and $\overline{CE}$}
\lqn{367}{Construct a triangle LMN with $m\overline{LM}=7\ cm,\ m\angle L=70^{\circ}$ and $m\angle M=45^{\circ}$ Find a point within the triangle LMN which is equidistant from L and M and 3 cm from L.}
\lqn{368}{Construct a right angled triangle RST with $m\overline{RS}=6.8\ cm,\ m\angle S=90^{\circ}$ and $m\overline{ST}=7.5\ cm$. Find a point within the triangle RST which is equidistant from $\overline{RS}$ and $\overline{RT}$ and 4.5 cm from R.}
\lqn{369}{Construct a rectangle UVWX with $m\overline{UV}=7.2\ cm$ and $m\overline{VW}=5.6cm$. Draw the locus of points at a distance of 2 cm from $\overline{UV}$ and 3.5 cm from W.}
\lqn{370}{Imagine two cell towers located at points A and B on a coordinate plane. The GPS-enabled device, positioned somewhere on the plane, receives signals form both towers. To ensure accurate navigation, the device is placed equidistant from both towers to estimate its position. Draw this locus of navigation.}
\lqn{371}{Epidemiologists use loci to determine infection zones, especially for contagious diseases, to predict the spread and take containment measures. In the case of a disease outbreak, authorities might determine a quarantine zone within 10 km of the infection sources. Draw the locus of all points 10 km from the source defining the quarantine area to monitor and control the disease's spread.}
\lqn{372}{There is a treasure buried somewhere on the island. The treasure is 24 kilometres from A and equidistant from B and C. Using a scale of 1cm to represent 10 km, find where the treasure could be buried.}
\lqn{373}{Thera is an apple tree at a distance of 90 metres from banana tree in the garden of Sara's house. Sara wants to plant a mango tree M which is 64 metres from apple tree and between 54 and 82 metres from the banana tree. Using a scale of 1cm to represent 10m, Find the points where the mango tree should be planted.}
\lqn{374}{Construct a triangle ABCD with $m\overline{AB}=5\ cm$ and $m\overline{BC}=3.2\ cm$. Draw the locus of all points which is: at a distance of 3.1 cm from point A. Label the point P inside the rectangle which is 3.1 cm from point A and equidistant from A and B.}
\lqn{375}{Construct a triangle ABCD with $m\overline{AB}=5\ cm$ and $m\overline{BC}=3.2\ cm$. Draw the locus of all points which is: equidistant from A and B. Label the point P inside the rectangle which is 3.1 cm from point A and equidistant from A and B.}
\lqn{376}{Construct an isosceles triangle DEF with vertical angle 80$^{o}$ at E and $m\overline{EF}=m\overline{DE}=4.8\ cm$. Draw the locus of all points which is: at a distance of 2.8 cm from point E. Label the point X inside the triangle which is 2.8 cm from point E and equidistant from $\overline{DE}$ and $\overline{EF}$.}
\lqn{377}{Construct an isosceles triangle DEF with vertical angle 80$^{o}$ at E and $m\overline{EF}=m\overline{DE}=4.8\ cm$. Draw the locus of all points which is: equidistant from $\overline{DE}$ and $\overline{EF}$. Label the point X inside the triangle which is 2.8 cm from point E and equidistant from $\overline{DE}$ and $\overline{EF}$.}
\lqn{378}{Construct a right angled triangle with measures of sides 6 cm, 8 cm and 10 cm.}
\lqn{379}{Construct a triangle ABC with $m\overline{AB}=5.3\ cm,\ m\angle A=30^{\circ}$ and $m\angle B=120^{\circ}$ Draw the locus of all points which are equidistant fro A and B.}
\lqn{380}{Construct a triangle with $m\overline{DE}=7.3\ cm,\ m\angle D=42^{\circ}$ and $m\overline{EF}=5.4\ cm$.}
\lqn{381}{Construct a triangle XYZ with $m\overline{YX}=8\ cm,\ m\overline{YZ}=7\ cm$ and $m\overline{XZ}=6.5\ cm$. Draw the locus of all points which are equidistant from $\overline{XY}$ and $\overline{XZ}$.}
\lqn{382}{Construct a triangle FGH such that $m\overline{FG}=m\overline{GH}=6.4\ cm,\ m\angle G=122^{\circ}$. Draw the locus of all points which is: equidistant from F and G.}
\lqn{383}{Construct a triangle FGH such that $m\overline{FG}=m\overline{GH}=6.4\ cm,\ m\angle G=122^{\circ}$. Draw the locus of all points which is: equidistant from $\overline{FG}$ and $\overline{GH}$.}
\lqn{384}{Construct a triangle FGH such that $m\overline{FG}=m\overline{GH}=6.4\ cm,\ m\angle G=122^{\circ}$. Draw the locus of all points which is: Mark the point where the two loci intersect.}
\lqn{385}{Two houses Q and R 73 metres apart. Using a scale of 1 cm to represent 10 m, construct the locus of a point P which moves such that it is: at a distance of 32 metres from Q.}
\lqn{386}{Two houses Q and R 73 metres apart. Using a scale of 1 cm to represent 10 m, construct the locus of a point P which moves such that it is: at a distance of 48 metres from the line joining Q and R.}
\lqn{387}{The field is in the form of a rectangle ABCD with $m\overline{AB}=70m$ and $m\overline{BC}=60\ m$. construct the rectangle ABCD using a scale of 1 cm to represent 10 m. Show the region inside the field which is less than 30 m from C and farther than 25 m from $\overline{AB}$.}
\lqn{388}{From the table given below. Draw a frequency polygon on histogram for the given frequency distribution:}
\dtab{|l|c|c|c|c|c|}{\textbf{Weight (kg)} & 10-14 & 15-19 & 20-24 & 25-29 & 30-34\\ \hline \textbf{Frequency (f)} & 06 & 17 & 23 & 30 & 22}
\lqn{389}{The marks obtained by the students of Grade-10 in mathematics test were grouped into the following frequency distribution:}
\dtab{l|c|c|c|c|c|c}{\textbf{Marks} & 35-37 & 38-44 & 45-54 & 55-61 & 62-67 & 68-72\\ \textbf{Frequency} & 1 & 2 & 16 & 13 & 9 & 3}
Draw a histogram for the above distribution.
\lqn{390}{Make a frequency polygon on histogram for the following grouped data:}
\dtab{l|c|c|c|c|c}{\textbf{Class limits} & 5-8 & 8-12 & 12-20 & 20-25 & 25-27\\ \textbf{Frequency (f)} & 2 & 12 & 25 & 32 & 14}
\lqn{391}{For a school staff, the following expenditure (rupee in hundred) are required for the repair of chairs:}
\nums{145\quad152\quad153\quad156\quad158\quad160\quad146\quad152\quad155\quad159\\ 161\quad163\quad165\quad147\quad148\quad151\quad154\quad156\quad158\quad160\\ 144\quad167\quad151\quad150\quad152\quad149\quad145\quad153\quad152\quad155}
prepare a frequency distribution by tally bar method using 3 as the size of class limits and write down what are the frequencies of the last three classes?
\lqn{392}{Given below are the weights in kg of 30 students of a high school:}
\nums{30\quad33\quad24\quad21\quad15\quad39\quad37\quad44\quad42\\ 33\quad28\quad29\quad32\quad31\quad28\quad26\quad32\quad34\\ 38\quad36\quad41\quad30\quad35\quad41\quad23\quad26\quad18}
Taking 5 as the size of the class limit, prepare a frequency table and construct a frequency polygon.
\lqn{393}{A group of Grade-10 student obtained the following marks out of 100 marks in English test:}
\nums{58\quad59\quad58\quad33\quad40\quad58\quad45\quad46\quad43\quad45\quad45\\ 50\quad52\quad49\quad50\quad57\quad52\quad55\quad49\quad50\quad62\quad49\\ 48\quad44\quad42\quad47\quad46\quad47\quad46\quad53\quad40\quad44}
Classify the data into a frequency distribution by (direct method) taking 6 as the size of class limit. Also find the class limit with least class frequency and construct histogram for the data.
\lqn{394}{Following are the number of telephone calls made in a week to 30 teachers of high school:}
\nums{5\quad8\quad11\quad25\quad13\quad16\quad20\quad17\quad15\quad16\quad30\quad21\quad14\quad18\\ 6\quad22\quad26\quad15\quad19\quad35\quad29\quad31\quad23\quad25\quad20\quad10\quad9\quad7}
Construct a frequency distribution with number of classes?
\lqn{395}{The frequency distribution of ages (in years) of 76 member of a locality is available. Draw a histogram for this data.}
\dtab{l|c|c|c|c|c|c|c}{\textbf{Class limits} & 2-4 & 4-9 & 9-12 & 12-17 & 17-20 & 20-27 & 27-30\\ \textbf{Frequency (f)} & 7 & 10 & 18 & 20 & 10 & 7 & 4}
\lqn{396}{The following are the marks obtained by 30 students out of 100 in the subject of Mathematics at their final examination. Construct frequency polygon for the following frequency table:}
\dtab{l|c|c|c|c|c}{\textbf{Marks} & 45-49 & 50-54 & 55-59 & 60-64 & 56-69\\ \textbf{Frequency} & 2 & 9 & 4 & 7 & 3}
\lqn{397}{The marks obtained by the students in the subject of English are given below:}
\dtab{l|c|c|c|c|c}{Marks obtained & 15-19 & 20-24 & 25-29 & 30-34 & 35-39\\ Frequency & 9 & 18 & 35 & 17 & 5}
Find: Arithmetic mean of their marks by direct and short formula.
\lqn{398}{The marks obtained by the students in the subject of English are given below:}
\dtab{l|c|c|c|c}{Marks obtained & 15-19 & 20-24 & 25-29 & 30-34\\ Frequency & 9 & 18 & 35 & 17}
Find: Median of their marks.
\lqn{399}{Given below is a frequency distribution:}
\dtab{l|c|c|c|c|c}{\textbf{Class interval} & 5-9 & 10-14 & 15-19 & 20-24 & 25-29\\ \textbf{Frequency} & 1 & 8 & 18 & 11 & 2}
Find the mode of the frequency distribution.
\lqn{400}{The monthly attendance of 10 students for their lunch in the hostel is recorded as: 21,15,16,18,14,17,15,12,13,11. Find the median and mode of the attendance. Also find the mean if D= A-20.}
\lqn{401}{On a prize distribution day, 50 students brought pocket money as under:}
\dtab{l|c|c|c|c|c}{\textbf{Rupees} & 5-10 & 10-15 & 15-20 & 20-25 & 25-30\\ \textbf{Frequency (f)} & 12 & 9 & 18 & 7 & 4}
Find the median and mode of the above data.
\lqn{402}{The children Haris, Maham and Minal made the following scores in a game conducted by a group of teachers in the school:}
\dtab{l|c|c|c|c}{\textbf{Haris scores} & 50 & 55 & 70 & 85\\ \textbf{Maham scores} & 75 & 60 & 60 & 45\\ \textbf{Minal scores} & 80 & 77 & 66 & 42}
It is decided that the candidate who gets the highest average score will be awarded rupees 1000. Who will get the awarded amount?
\lqn{403}{Being partners Hafsa and Fatima took part in a quiz programme. They made the following numbers of points 45,51,58,61,74,48,46 and 50. Compute the average number of points using deviation D=x=-58.}
\lqn{404}{A person purchased the following food items:}
\dtab{|c|c|c|}{\textbf{Food item} & \textbf{Quantity (in kg)} & \textbf{Cost per kg (in Rs.)}\\ \hline Rice & 10 & 96\\ Flour & 12 & 48\\ Ghee & 4 & 190\\ Sugar & 3 & 49\\ Mutton & 2 & 650}
What is the weighted mean of cost of food items per kg?
\lqn{405}{The arithmetic mean of the ages of 20 boys is 13 years, 4 months and 5 days. Find the sum of their ages. If one of the boys is age exactly 15 years. What is the average age of the remaining boys?}
\lqn{406}{Given below is the frequency distribution derived by making a substitution as D=X-20. Calculate the arithmetic mean:}
\dtab{c|c|c|c|c|c|c}{\textbf{D} & -6 & -4 & -2 & 0 & 2 & 4\\ \textbf{f} & 1 & 3 & 6 & 20 & 26 & 12}
\lqn{407}{The marks of Mathematics of Jamal in eight monthly tests were 75,76,80,80,82,82,82,85. Find the mode of the marks.}
\lqn{408}{Ten students were asked about the number of questions they have solved out of 20 questions last week. records were 13,14,15,11,16,10,19,20,18,17. Find the mode of the data.}
\lqn{409}{A survey was conducted from the 15 students of a school and asked the students about their favourite colour. The response are: Purple, yellow, purple, yellow,yellow, red,blue,green, yellow,yellow, red, blue, yellow, purple, green. Find mode of the data.}
\lqn{410}{The student each from Grade V section A and B of a well reputed school. were taken randomly. Their weights were measured in kg. and recorded as given below:}
\dtab{|l|c|c|c|c|c|c|c|c|c|}{\textbf{Weights (kg) Section A} & 30 & 28 & 32 & 29.5 & 35 & 34 & 31 & 33 & 40\\ \textbf{Weights (kg) Section B} & 35 & 31.5 & 34.5 & 35 & 32.8 & 38 & 39.5 & 36 & 36.}
\lqn{411}{A government allocates funds of Rs, 200,000 to five sectors of a school: (i) school library: Rs, 35,000 (ii) Sports facilities: Rs. 25,000 (iii) Parking Area: Rs, 40,000 (iv) Room renovation: Rs, 45,000 (v) Furniture: Rs, 55,000. Find the average of the fund allocation in each sector of a school.}
\lqn{412}{Given below are the marks out of 100 obtained by 100 students in a examination. Find the average marks of the students:}
\dtab{c|c|c|c|c|c}{\textbf{Marks} & 30-35 & 35-40 & 40-45 & 45-50 & 50-55\\ \textbf{No. of students} & 14 & 16 & 18 & 23 & 18}
\lqn{413}{Deviations from 12.5 of ten different values are 6,-2,3.5,9,8.7,-5.5,14,11.3,-6.8,-4.2, find the arithmetic mean.}
\lqn{414}{The heights (in inches) of 200 students are recorded in the following frequency distribution. Find the mean height of the student by short formula:}
\dtab{|c|c|c|c|c|c|c|c|c|}{\textbf{Height (x) (in inches)} & 51 & 52 & 53 & 54 & 55 & 56 & 57 & 58\\ \textbf{Frequency (f)} & 2 & 5 & 8 & 24 & 55 & 45 & 38 & 16}
\lqn{415}{The heights of 100 athletes, measured to the nearest (inches) are given in the following table. Find the median:}
{\scriptsize\setlength{\tabcolsep}{3pt}\dtab{|c|c|c|c|c|c|c|c|c|c|}{\textbf{Heights (in inches)} & 62.5-63.5 & 63.5-64.5 & 64.5-65.5 & 65.5-66.5 & 66.5-67.5 & 67.5-68.5 & 68.5-69.5 & 69.5-70.5 & 70.5-71.5\\ \textbf{No. of Students} & 4 & 6 & 10 & 20 & 30 & 13 & 12 & 3 & 2}}
\lqn{416}{Following are the weights (in kg) of 50 men. Find the median weight.}
\dtab{|c|c|c|c|c|c|}{\textbf{Weights (kg)} & 110-114 & 115-119 & 120-124 & 125-129 & 13\\ \textbf{No. of men (f)} & 5 & 12 & 23 & 6 &}
\lqn{417}{Following are the heights in (inches) of 40 students in Grade 8.}
\dtab{|c|c|c|c|c|c|c|}{Heights (inches) & 48-50 & 50-52 & 52-54 & 54-56 & 56-58 & 58-60\\ No. of students (f) & 5 & 7 & 10 & 9 & 6 & 3}
find mode of above data.
\lqn{418}{A medicine company started marketing of a sample for medicine in seven different areas of a city. The company distributed the packets of medicine in each area of the city and the weight of each area based on the demand of the medicine. Find the mean and weighted mean of the given data.}
\dtab{|c|c|c|}{\textbf{Area} & \textbf{Number of packets} & \textbf{Weights (Kg)}\\ \hline A & 15 & 5\\ B & 25 & 4\\ C & 18 & 3\\ D & 23 & 4\\ E & 15 & 2\\ F & 10 & 1\\ G & 8 & 2}
\lqn{419}{A toy factory sold toys in a month. Consider the following data:}
\dtab{|c|c|c|c|c|c|}{Class limits & 10-20 & 20-30 & 30-40 & 40-50 & 50-60\\ f & 15 & 28 & 45 & 29 & 20}
(i)Calculate mean, median and mode for the number of toys sold by the factory.\par
\hspace*{3em}(ii) Also tell the modal class of the distribution.
\lqn{420}{Following are the weights of 40 students recorded to the nearest (lbs).}
\nums{138\quad164\quad150\quad132\quad144\quad125\quad149\quad157\quad146\\ 140\quad147\quad136\quad148\quad152\quad144\quad168\quad126\quad138\\ 163\quad119\quad154\quad165\quad146\quad173\quad142\quad147\quad135\\ 140\quad135\quad161\quad145\quad135\quad142\quad150\quad156\quad145}
make a frequency table taking size of class limit as 10. Also draw a frequency polygon of the given data.
\lqn{421}{Adjoining distribution showed maximum load (in kg) supported by certain ropes. Find the mean load using short method:}
\dtab{c|c|c|c|c|c|c}{\textbf{Max-Load kg} & 93-97 & 98-102 & 103-107 & 108-112 & 113-117 & 118-122\\ \textbf{No. of ropes} & 2 & 5 & 8 & 12 & 6 & 2}
\lqn{422}{Given below is frequency distribution: Draw frequency polygon and histogram for the distribution:}
\dtab{c|c|c|c|c|c}{Class limits & 5-9 & 10-14 & 15-19 & 20-24 & 25-29\\ Frequency & 1 & 8 & 18 & 11 & 2}
\lqn{423}{For the following data, find the weighted mean:}
\dtab{|c|c|c|}{\textbf{Item} & \textbf{Quantity} & \textbf{Cost of item (Rs.)}\\ \hline Chair & 20 & 500\\ Table & 20 & 400\\ Black board & 10 & 750\\ Tube light & 25 & 230\\ Cupboard & 09 & 950}
\lqn{424}{Two partners Mr. Aslam and Mrs. Kalsoom run a company. In the following data the weekly wages (in Rs) of employees who work in the company are given:}
\dtab{|c|c|c|c|c|c|}{\textbf{Wages (Rs.)} & 600-700 & 700-800 & 800-900 & 900-1000 & 1000\\ \textbf{Employees} & 3 & 5 & 7 & 21 &}
Find the arithmetic mean, median, and standard deviation of the given table.
\lqn{425}{From the table given below. Draw a frequency polygon on histogram for the given frequency distribution:}
\dtab{|c|c|c|c|c|c|c|}{\textbf{Weight (kg)} & 50-56 & 57-59 & 60-64 & 65-72 & 73-75 & 76-80\\ \textbf{Frequency (f)} & 25 & 32 & 40 & 30 & 15 & 8}
\lqn{426}{A principal of a school allocates funds of Rs. 50,000 to five different sectors: (i) Chairs Rs, 1500, (ii) Tables Rs. 12,000 (iii) Black boards: Rs. 6,000 (iv) Room renovation: Rs. 10,000 (v) Gardening: Rs. 7,000. Find the average of funds allocation in each sector of the school.}
\lqn{427}{Aslam rolled a dice. What is the probability of getting the number 3 or 4? Also find the probability of not getting the number 3 or 4}
\lqn{428}{Abdul Hadi labelled cards from 1 to 30 and put them in a box. He selects a card at random. What is the probability that selected card containing:\\
\hspace*{0pt}(i) the number 25 (ii) number between 17 to 22\\
(iii) number at least 20 (iv) number not 27 and 29\\
(v) number not between 12 to 15}
\lqn{429}{The probability that Ayesha will pass the examination is 0.85. What.will be the probability that Ayesha will not pass the examination?}
\lqn{430}{A card is selected at random from well shuffled pack of 52 plying cards. What will be the probability of selecting:\\
(i) a queen (ii) neither a queen nor a jack}
\lqn{431}{A card is chosen at random from a pack of 52 playing cards. Find the probability of getting: (i) a jack (ii) no diamond}
\lqn{432}{Taabish tossed a fair coin and rolled a fair dice once. Find the probability of the following events:\\
(i) tail on coin and at least 4 on dice.\\
(ii) head on coin and the number 2,3 on dice.\\
(iii) head and tail on coin and the number 6 on dice.\\
(iv) not tail on coin and the number 5 on dice.\\
(v) not head on coin and the number 5 and 2 on dice.}
\lqn{433}{Let A,B and C are the missiles and they are fired at a target, if the probabilities of hitting the target,\\
$P(A)=\dfrac{1}{4},\ P(B)=\dfrac{3}{7},\ P(C)=\dfrac{5}{9}$, respectively. Find the probabilities of\\
(i) missile A does not hit the target.\\
(ii)missile B does not hit the target.\\
(iii)missile C does not hit the target.}
\lqn{434}{A ball contains 5 blue balls and 8 green balls. Find the probability of selecting at random:\\
(i) a blue ball (ii) a green ball (iii) not a green ball}
\lqn{435}{A researcher collected data on number of deaths from Horse-Ricks in Russian Army crops over to years. Thetable is as follows:}
\dtab{|c|c|c|c|c|c|c|}{\textbf{No.of Death} & 0 & 1 & 2 & 3 & 4 & 5\\ \textbf{Frequency} & 60 & 50 & 87 & 40 & 32 & 15}
Find the relative frequency of the given data.
\lqn{436}{The frequency of defective products in 750 samples are shown in the following table. Find the relative frequency for the given table.}
\dtab{|c|c|c|c|c|c|c|c|}{\textbf{No.of Defective per sample} & 0 & 1 & 2 & 3 & 4 & 5 & 6\\ \textbf{No. of Sample} & 120 & 40 & 94 & 85 & 105 & 50 & 40}
\lqn{437}{A quiz competition on general knowledge is conducted. The number of corrected answers out of 5 questions for 100 sets of questions is given below.}
\dtab{|c|c|c|c|c|c|c|}{\textbf{x} & 0 & 1 & 2 & 3 & 4 & 5\\ \textbf{f} & 10 & 23 & 15 & 25 & 18 & 9}
Find the relative frequencies for the given data.
\lqn{438}{A survey was conducted from the 108 students of a class and asked about their favourite food. The responses are as under:}
\dtab{|c|c|c|c|c|c|}{\textbf{Name of food item} & Biryani & Fresh Juice & Chicken & Bar.B.Q & S\\ \textbf{No. of students} & 40 & 07 & 21 & 15 &}
(i) how many percentages of students like biryani?\par
\hspace*{3em}(ii) how many percentages of students like chicken`\par
\hspace*{3em}(iii) which food is the least like by the students?\par
\hspace*{3em}(iv) which food is the most prefer by the students?
\lqn{439}{Find the expected frequencies of the givenjdata if the experiment is repeated 200 times.}
\dtab{|c|c|c|c|c|c|c|c|}{\textbf{x} & 0 & 1 & 2 & 3 & 4 & 5 & 6\\ \textbf{P(x)} & 0.11 & 0.21 & 0.17 & 0.18 & 0.09 & 0.17 & 0.07}
\lqn{440}{The probability of getting`S sixes while tossing six dice is $\dfrac{2}{5}$, the dice is rolled 200 times. Howanany times would you expect it to show 5 sixes?}
\lqn{441}{Find the relative frequency of the given data.}
\dtab{|c|c|c|c|c|c|c|c|}{\textbf{x} & 2 & 3 & 4 & 5 & 6 & 7 & 8\\ \textbf{f} & 3 & 5 & 6 & 9 & 10 & 8 & 2}
\lqn{442}{Abdul Rehman obtained different marks in different subjects out of 100 marks. The detail is as under:}
\dtab{|c|c|c|c|c|c|c|}{\textbf{Subjects} & English & urdu & islamiyat & Mathematics & Science & Computer Science\\ \textbf{marks obtained} & 75 & 80 & 72 & 95 & 81 & 85}
Find the relative frequency of above given data.
\lqn{443}{Six fair dice are rolled 50 times. The probability of occurrence of different number of sixes are given below. Find the expected frequency of the following data:}
\dtab{|c|c|c|c|c|c|c|c|}{\textbf{x} & 0 & 1 & 2 & 3 & 4 & 5 & 6\\ \textbf{P(x)} & 0.09 & 0.10 & 0.12 & 0.24 & 0.10 & 0.20 & 0.15}
Find the expected frequency of occurrence of each six.
\lqn{444}{An urn contains 10 red balls, 5 green balls and 8 blue balls. Find the probability of selecting at random.\\
(i) a green ball (ii) a red ball (iii) a blue ball (iv) nota red ball (v) not a green ball}
\lqn{445}{Three coins are tossed together. what is the probability of getting:\\
(i) exactly three heads (ii) at least two tails\\
(iii) not at least two heads (iv) not exactly two heads}
\lqn{446}{A card is drawn from a well shuffled pack of 52 playing cards. What will be the probability of getting:\\
(i) king or jack of red colour (ii) not 2 of club and spade}
\lqn{447}{Six coins are tossed 600 times. The number of occurrences of tails are recorded and shown in the table given below:}
\dtab{|c|c|c|c|c|c|c|c|}{\textbf{No. of Tails} & 0 & 1 & 2 & 3 & 4 & 5 & 6\\ \textbf{Frequency} & 110 & 90 & 105 & 80 & 76 & 123 & 16}
Find the relative frequency of given table:
\lqn{448}{From a lot containing 25 items, 8 items are defective. Find the relative frequency of non-defective items, also find the expected frequency of non-defective items.}

\end{document}
